\section{Большой нормальный градиент и неизменная траектория}
\label{sec:normal}
Правая часть условия гладкости может расти из-за линейного слагаемого,
которое не меняет кривизну. Следующее утверждение описывает семейство,
для которого такой рост не влияет на метод с подбором шага.

\begin{proposition}[Неизменность метода при нормальном линейном слагаемом]
\label{prop:normal_invariance}
Пусть $Q_0$ компактно и выпукло, а $\varphi$ выпукла и принадлежит
$C^2$ в окрестности $Q_0$. Для $A\ge0$ рассмотрим
\[
 Q=Q_0\times[0,1],\qquad f_A(z,y)=\varphi(z)+Ay.
\]
Запустим точный метод раздела~\ref{sec:stopping} из $(z_0,0)$ с
одинаковыми точностью и начальной кривизной. При неоднозначности
выбираем один и тот же линейный минимум на $Q_0$, а по координате $y$
выбираем ноль. Тогда для всех $A\ge0$ совпадают точки, зазоры,
принятые шаги, пробы подбора шага и числа обращений к оракулам.
Если $\varphi$ удовлетворяет~\eqref{eq:gradient_smoothness}, то
каждая $f_A$ удовлетворяет ему с теми же константами. В минимуме
$(z^*,0)$ имеем
\[
 \|\nabla f_A(z^*,0)\|
   =\sqrt{\|\nabla\varphi(z^*)\|^2+A^2}\ge A.
\]
\end{proposition}
\begin{proof}
В точке $(z,0)$ градиент равен $(\nabla\varphi(z),A)$. Поэтому полный
линейный оракул возвращает $(s,0)$, где $s$ минимизирует
$\langle\nabla\varphi(z),s\rangle$ на $Q_0$.
Последняя координата направления равна нулю. Его норма, направленная
производная, зазор и все значения функции вдоль него не зависят от $A$.
Значит, совпадают сохраняемая нижняя граница и все проверки шага.
Индукция доказывает совпадение всего выполнения алгоритма.
Гессиан состоит из блока $\nabla^2\varphi$ и нулевого блока, а
$\|\nabla f_A\|\ge\|\nabla\varphi\|$, что доказывает утверждение о
гладкости.
\end{proof}

Верхняя граница $\|\nabla^2\varphi\|$ на $Q_0$ даёт значение $K_Q$,
не зависящее от $A$, для теоремы~\ref{thm:finite_stopping}.
Таким образом, и выполнение метода, и верная оценка числа шагов
могут не зависеть от этой нормальной компоненты.
Обратная матрица Гессе нигде не используется.

\subsection{Сравнение с шагом по фиксированным параметрам гладкости}
Правило из Algorithm~1 работы~\cite{vyguzov} имеет вид
\begin{equation}\label{eq:vyguzov_step}
 \alpha_k=\min\left\{1,
 \frac{G(x_k)}{e(L_0+L_1\|\nabla f(x_k)\|)\|s_k-x_k\|^2}
 \right\}.
\end{equation}
Здесь $e$ является основанием натурального логарифма. Сравниваем с
подбором шага именно это правило. Адаптивный Algorithm~2 из той же
работы является другим методом.

\begin{theorem}[Линейное различие по нормальному коэффициенту]
\label{thm:normal_separation}
Пусть
\[
 f_A(x,y)=e^x+Ay,\qquad Q=[-1,1]\times[0,1],\qquad A>0.
\]
Начальная точка равна $(a,0)$, где $0<a<1$, и
$0<\varepsilon<1-e^{-1}$. В правиле~\eqref{eq:vyguzov_step} положим
$L_0=L_1=1$ и используем тот же критерий остановки с сохраняемой
нижней границей, что и в разделе~\ref{sec:stopping}. До остановки
этому правилу требуется не менее
\begin{equation}\label{eq:normal_lower}
 \left\lceil a e^{1-a}A\right\rceil
\end{equation}
шагов. Выполнение метода с подбором шага, включая конечный момент
остановки, не зависит от $A$.
\end{theorem}
\begin{proof}
Функции выпуклы на всей плоскости, причём
\[
 \nabla f_A(x,y)=(e^x,A),\qquad
 \nabla^2f_A(x,y)=\operatorname{diag}(e^x,0).
\]
Поэтому условие~\eqref{eq:gradient_smoothness} выполнено всюду с
одними и теми же $L_0=L_1=1$. Минимум на $Q$ находится в $(-1,0)$.
Полный линейный оракул всегда возвращает $(-1,0)$. Точки остаются
на $y=0$, а первая координата убывает. В текущей координате $x$
зазор равен $e^x(x+1)$. Правило~\eqref{eq:vyguzov_step} уменьшает
координату не более чем на
\[
 \frac{e^x}{e(1+\sqrt{e^{2x}+A^2})}
 \le\frac{e^a}{eA}.
\]
Оценка остаётся верной и при ограничении шага единицей. Пока
координата не стала отрицательной, истинная ошибка не меньше
$1-e^{-1}>\varepsilon$, поэтому верный сертификат не допускает
остановку. Требуется пройти расстояние не меньше $a$, откуда
следует~\eqref{eq:normal_lower}.
Независимость подбора шага от $A$ доказана в
утверждении~\ref{prop:normal_invariance}. Его конечная остановка
следует из теоремы~\ref{thm:finite_stopping}, например при $K_Q=e$
и $D=\sqrt5$. Обе константы не зависят от $A$.
\end{proof}

Результат относится к конкретному выпуклому семейству и указанному
правилу шага. Причина видна в~\eqref{eq:vyguzov_step}: нормальная
компонента входит в знаменатель, хотя не меняет ни функцию вдоль
сделанных шагов, ни гессиан. Подбор шага проверяет фактическое
изменение функции вдоль направления.
